% Vietnamese translation for OI Public Library
% Source-PDF: IOI中国国家候选队论文/2004/杨思雨.pdf
\documentclass[11pt]{article}
\usepackage{oipl-vn}
\usepackage{graphicx}
\renewcommand{\figurename}{Hình}

\title{Các thao tác cơ bản và ứng dụng của cây splay}
\author{Yang Siyu}
\date{}

\begin{document}
\maketitle

\origtitle{Basic operations and applications of splay trees}
\sourcepath{IOI China candidate team papers / 2004 / Yang Siyu.pdf}

\section*{Từ khóa}

Cây splay; thao tác cơ bản; ứng dụng.

\section*{Tóm tắt}

Bài viết này giới thiệu các thao tác cơ bản của cây splay và cách dùng nó
trong giải bài tập. Nội dung chính có thể chia thành bốn phần.

Phần mở đầu nêu vai trò quan trọng của cây tìm kiếm nhị phân trong thi lập
trình, đồng thời chỉ ra rằng trong một số trường hợp độ phức tạp thời gian
của cây tìm kiếm nhị phân có thể rất cao. Từ đó ta dẫn tới cây splay, một
biến thể có độ phức tạp trung bình theo nghĩa khấu hao tốt hơn.

Phần thứ hai trình bày các thao tác cơ bản của cây splay, kèm phân tích và
chứng minh độ phức tạp thời gian. Kết quả là mọi thao tác cơ bản của cây
splay đều có độ phức tạp khấu hao $O(\log n)$; vì vậy cây splay có thể xem
là một loại cây tìm kiếm nhị phân cân bằng theo nghĩa khấu hao.

Phần thứ ba dùng một ví dụ để minh họa ứng dụng của cây splay trong giải
bài, và so sánh nó với một số cấu trúc dữ liệu dạng cây khác.

Phần cuối cùng tổng kết các ưu điểm của cây splay.

\section{Mở đầu}

Cây tìm kiếm nhị phân có thể hỗ trợ nhiều thao tác trên tập động. Vì vậy,
trong thi tin học, cây tìm kiếm nhị phân giữ một vị trí rất quan trọng: nó
có thể dùng để biểu diễn tập có thứ tự, xây dựng chỉ mục, cài đặt hàng đợi
ưu tiên, v.v.

Thời gian của các thao tác cơ bản trên cây tìm kiếm nhị phân tỉ lệ với chiều
cao của cây. Với một cây nhị phân hoàn chỉnh có $n$ nút, thời gian xấu nhất
của những thao tác này là $O(\log n)$. Nhưng nếu cây có $n$ nút lại suy biến
thành một dây chuyền, thời gian xấu nhất sẽ là $O(n)$. Một số biến thể của
cây tìm kiếm nhị phân, như cây đỏ đen hay cây AVL, vẫn bảo đảm hiệu năng tốt
trong trường hợp xấu nhất.

Cây splay cũng là một cải tiến của cây tìm kiếm nhị phân. Nó không bảo đảm
rằng cây luôn ``cân bằng'', nhưng với một dãy thao tác trên cây splay, ta có
thể chứng minh chi phí khấu hao của mỗi thao tác đều là $O(\log n)$. Theo
nghĩa đó, cây splay cũng là một cây tìm kiếm nhị phân cân bằng. Trong các cấu
trúc dữ liệu dạng cây, cây splay còn có ưu điểm tốt về yêu cầu bộ nhớ và độ
phức tạp khi lập trình.

\section{Các thao tác cơ bản của cây splay}

Cây splay là một cải tiến của cây tìm kiếm nhị phân, nên nó cũng giữ tính thứ
tự: với mỗi nút $x$, mọi phần tử trong cây con trái của $x$ đều nhỏ hơn $x$,
và mọi phần tử trong cây con phải của $x$ đều lớn hơn $x$. Điểm khác với cây
tìm kiếm nhị phân thông thường là cây splay có khả năng tự điều chỉnh, nhờ
thao tác splay $\operatorname{Splay}(x,S)$.

\subsection{Thao tác splay}

Thao tác $\operatorname{Splay}(x,S)$ đưa phần tử $x$ trong cây splay $S$ lên
gốc bằng một dãy phép quay, trong khi vẫn giữ nguyên tính thứ tự của cây.
Trong quá trình điều chỉnh, ta xét ba trường hợp.

\paragraph{Trường hợp 1.}
Cha $y$ của $x$ là gốc. Nếu $x$ là con trái của $y$, ta thực hiện một phép
\term{Zig}, tức quay phải tại $y$; nếu $x$ là con phải của $y$, ta thực hiện
một phép \term{Zag}, tức quay trái tại $y$. Sau phép quay, $x$ trở thành gốc
của cây $S$ và thao tác kết thúc.

\begin{figure}[ht]
\centering
\includegraphics[width=0.62\linewidth]{assets/ioi/2004/splay-yang-siyu/fig-000.png}
\caption{Trường hợp Zig/Zag: cha của $x$ là gốc, nên một phép quay đơn đưa
$x$ lên gốc.}
\end{figure}

\paragraph{Trường hợp 2.}
Cha $y$ của $x$ không phải là gốc; gọi $z$ là cha của $y$. Nếu $x$ và $y$
cùng là con trái của cha mình, hoặc cùng là con phải của cha mình, ta thực
hiện một bước \term{Zig-Zig} hoặc \term{Zag-Zag}. Nghĩa là ta quay ở nút
$z$, rồi tiếp tục quay ở nút $y$ theo cùng một hướng.

\begin{figure}[ht]
\centering
\includegraphics[width=0.62\linewidth]{assets/ioi/2004/splay-yang-siyu/fig-001.png}
\caption{Trường hợp Zig-Zig/Zag-Zag: $x$ và $y$ nằm cùng phía so với cha của
mình, nên thực hiện hai phép quay cùng hướng.}
\end{figure}

\paragraph{Trường hợp 3.}
Cha $y$ của $x$ không phải là gốc; gọi $z$ là cha của $y$. Nếu một trong hai
nút $x,y$ là con trái của cha mình còn nút kia là con phải, ta thực hiện
\term{Zig-Zag} hoặc \term{Zag-Zig}. Nghĩa là ta quay ở $y$ trước, rồi quay ở
$z$ theo hướng ngược lại.

\begin{figure}[ht]
\centering
\includegraphics[width=0.70\linewidth]{assets/ioi/2004/splay-yang-siyu/fig-002.png}
\caption{Trường hợp Zig-Zag/Zag-Zig: $x$ và $y$ nằm khác phía, nên hai phép
quay có hướng ngược nhau.}
\end{figure}

Ví dụ, khi thực hiện $\operatorname{Splay}(1,S)$, phần tử $1$ được đưa lên
gốc của cây splay $S$. Nếu sau đó thực hiện $\operatorname{Splay}(2,S)$, cây
thường trở nên cân bằng hơn theo trực giác. Bản thân quá trình splay không
phức tạp: ở mỗi bước chỉ cần phân loại vị trí của nút hiện tại rồi thực hiện
các phép quay tương ứng. Ba kiểu quay trên đều được tạo từ hai phép quay cơ
bản là quay trái và quay phải, nên việc cài đặt khá đơn giản.

\begin{figure}[ht]
\centering
\includegraphics[width=0.95\linewidth]{assets/ioi/2004/splay-yang-siyu/fig-003.png}
\caption{Ví dụ thực hiện $\operatorname{Splay}(1,S)$: các bước Zig-Zig rồi
Zig đưa nút $1$ lên gốc.}
\end{figure}

\begin{figure}[ht]
\centering
\includegraphics[width=0.82\linewidth]{assets/ioi/2004/splay-yang-siyu/fig-004.png}
\caption{Tiếp tục thực hiện $\operatorname{Splay}(2,S)$ sau ví dụ trước:
cây trở nên cân bằng hơn theo trực giác.}
\end{figure}

\subsection{Các thao tác cơ bản}

Dựa vào thao tác splay, ta có thể thực hiện các phép toán sau trên cây splay
$S$.

\paragraph{\texorpdfstring{$\operatorname{Find}(x,S)$}{Find(x,S)}.}
Kiểm tra xem phần tử $x$ có thuộc tập có thứ tự được biểu diễn bởi $S$ hay
không. Trước hết ta tìm $x$ giống như trên cây tìm kiếm nhị phân. Nếu $x$ có
trong cây, ta tiếp tục thực hiện $\operatorname{Splay}(x,S)$ để điều chỉnh
cây.

\paragraph{\texorpdfstring{$\operatorname{Insert}(x,S)$}{Insert(x,S)}.}
Chèn phần tử $x$ vào tập có thứ tự được biểu diễn bởi $S$. Trước hết ta chèn
$x$ vào đúng vị trí như trong cây tìm kiếm nhị phân thông thường, sau đó thực
hiện $\operatorname{Splay}(x,S)$.

\paragraph{\texorpdfstring{$\operatorname{Delete}(x,S)$}{Delete(x,S)}.}
Xóa phần tử $x$ khỏi tập có thứ tự được biểu diễn bởi $S$. Trước hết ta tìm
vị trí của $x$ bằng cách tìm kiếm trong cây nhị phân. Nếu $x$ không có con
hoặc chỉ có một con, ta xóa trực tiếp $x$, rồi dùng thao tác splay đưa cha
của $x$ lên gốc. Nếu $x$ có hai con, ta tìm nút kế tiếp $y$ của $x$, dùng
$y$ thay cho vị trí của $x$, rồi thực hiện $\operatorname{Splay}(y,S)$ để đưa
$y$ lên gốc.

\paragraph{\texorpdfstring{$\operatorname{Join}(S_1,S_2)$}{Join(S1,S2)}.}
Hợp nhất hai cây splay $S_1$ và $S_2$ thành một cây splay, với điều kiện mọi
phần tử của $S_1$ đều nhỏ hơn mọi phần tử của $S_2$. Ta tìm phần tử lớn nhất
$x$ trong $S_1$, dùng $\operatorname{Splay}(x,S_1)$ để đưa $x$ lên gốc của
$S_1$, rồi gắn $S_2$ làm cây con phải của $x$. Kết quả thu được là cây splay
mới $S$.

\begin{figure}[ht]
\centering
\includegraphics[width=0.74\linewidth]{assets/ioi/2004/splay-yang-siyu/fig-005.png}
\caption{Thao tác Join: đưa phần tử lớn nhất của $S_1$ lên gốc rồi gắn
$S_2$ làm cây con phải.}
\end{figure}

\paragraph{\texorpdfstring{$\operatorname{Split}(x,S)$}{Split(x,S)}.}
Dùng $x$ làm mốc để tách $S$ thành hai cây splay $S_1$ và $S_2$, trong đó mọi
phần tử của $S_1$ nhỏ hơn $x$, còn mọi phần tử của $S_2$ lớn hơn $x$. Trước
hết ta thực hiện $\operatorname{Find}(x,S)$ để đưa $x$ lên gốc của cây. Khi
đó cây con trái của $x$ chính là $S_1$, còn cây con phải là $S_2$.

\begin{figure}[ht]
\centering
\includegraphics[width=0.74\linewidth]{assets/ioi/2004/splay-yang-siyu/fig-006.png}
\caption{Thao tác Split: đưa $x$ lên gốc rồi tách hai cây con trái và phải
thành $S_1$ và $S_2$.}
\end{figure}

Ngoài năm thao tác trên, cây splay còn hỗ trợ tìm giá trị lớn nhất, nhỏ nhất,
tiền nhiệm, kế nhiệm, v.v. Những thao tác này cũng đều xây dựng trên nền tảng
của thao tác splay.

\subsection{Phân tích độ phức tạp thời gian}

Từ cách cài đặt các thao tác trên, có thể thấy hiệu quả thời gian của chúng
phụ thuộc hoàn toàn vào độ phức tạp của thao tác splay. Ta sẽ dùng phương
pháp kế toán để phân tích chi phí khấu hao của splay.

Trước hết định nghĩa một số ký hiệu. $S(x)$ là cây con gốc $x$; $|S|$ là số
nút của cây splay $S$. Đặt
\[
  \mu(S)=\lfloor \log |S| \rfloor,\qquad
  \mu(x)=\mu(S(x)).
\]

\begin{figure}[ht]
\centering
\includegraphics[width=0.52\linewidth]{assets/ioi/2004/splay-yang-siyu/fig-007.png}
\caption{Ví dụ ký hiệu $S(x)$ và các giá trị $\mu(x)$ trong phân tích khấu
hao của thao tác splay.}
\end{figure}

Ta dùng $1$ đồng để biểu diễn một đơn vị chi phí, trong đó một lần truy cập
vào một nút hoặc một phép quay được xem là chi phí đơn vị. Bất biến của cây
splay được định nghĩa như sau: tại mọi thời điểm, mỗi nút $x$ trong cây phải
có ít nhất $\mu(x)$ đồng tiền dự trữ.

Trong quá trình splay, chi phí được dùng cho hai việc:
\begin{enumerate}
  \item Trả cho thời gian sử dụng, tức mỗi thao tác đơn vị cần trả $1$ đồng.
  \item Khi hình dạng cây thay đổi, ta phải nạp thêm tiền hoặc phân phối lại
    số tiền đang có ở các nút để bất biến tiếp tục đúng.
\end{enumerate}

Ta có định lý sau về chi phí của thao tác splay.

\paragraph{Định lý.}
Trong mỗi thao tác $\operatorname{Splay}(x,S)$, tổng chi phí để điều chỉnh
cấu trúc cây và duy trì bất biến không vượt quá $3\mu(S)+1$.

\paragraph{Chứng minh.}
Ký hiệu $\mu(x)$ và $\mu'(x)$ lần lượt là lượng tiền tại nút $x$ trước và sau
một bước Zig, Zig-Zig hoặc Zig-Zag. Ta xét ba trường hợp.

\paragraph{Trường hợp 1: Zig hoặc Zag.}
Để duy trì bất biến của cây splay sau phép quay, chi phí thay đổi tiền dự
trữ là

\begin{figure}[ht]
\centering
\includegraphics[width=0.62\linewidth]{assets/ioi/2004/splay-yang-siyu/fig-008.png}
\caption{Ký hiệu trước và sau một bước Zig trong chứng minh chi phí khấu hao.}
\end{figure}

\[
\begin{aligned}
  \mu'(x)+\mu'(y)-\mu(x)-\mu(y)
    &= \mu'(y)-\mu(x) \\
    &\le \mu'(x)-\mu(x) \\
    &\le 3(\mu'(x)-\mu(x)) \\
    &= 3(\mu(S)-\mu(x)).
\end{aligned}
\]
Ngoài ra ta cần thêm $1$ đồng để trả cho thao tác truy cập và quay. Vì vậy
một bước Zig hoặc Zag tốn nhiều nhất $3(\mu(S)-\mu(x))+1$.

\paragraph{Trường hợp 2: Zig-Zig hoặc Zag-Zag.}
Để duy trì bất biến, chi phí thay đổi tiền dự trữ là

\begin{figure}[ht]
\centering
\includegraphics[width=0.62\linewidth]{assets/ioi/2004/splay-yang-siyu/fig-009.png}
\caption{Ký hiệu trước và sau một bước Zig-Zig trong chứng minh chi phí khấu hao.}
\end{figure}

\[
\begin{aligned}
  &\mu'(x)+\mu'(y)+\mu'(z)-\mu(x)-\mu(y)-\mu(z) \\
  &\quad= \mu'(y)+\mu'(z)-\mu(x)-\mu(y) \\
  &\quad= (\mu'(y)-\mu(x))+(\mu'(z)-\mu(y)) \\
  &\quad\le (\mu'(x)-\mu(x))+(\mu'(x)-\mu(x)) \\
  &\quad= 2(\mu'(x)-\mu(x)).
\end{aligned}
\]
Giống trường hợp trước, ta còn phải trả $1$ đồng cho thao tác đơn vị. Nếu
$\mu'(x)>\mu(x)$, thì
\[
  2(\mu'(x)-\mu(x))+1 \le 3(\mu'(x)-\mu(x)).
\]
Vì vậy chi phí của bước Zig-Zig không vượt quá $3(\mu'(x)-\mu(x))$.

Nếu $\mu'(x)=\mu(x)$, ta chứng minh được
\[
  \mu'(x)+\mu'(y)+\mu'(z)
  < \mu(x)+\mu(y)+\mu(z).
\]
Nghĩa là không cần nạp thêm tiền để duy trì bất biến, thậm chí còn thu hồi
được tiền; dùng $1$ đồng trong phần thu hồi này để trả cho thao tác truy cập
và quay. Để thấy điều đó, giả sử ngược lại rằng vế trái lớn hơn vế phải. Khi
đó ta có $\mu(x)=\mu'(x)=\mu(z)$, suy ra
$\mu(x)=\mu(y)=\mu(z)$ và $\mu(x)=\mu'(z)=\mu(z)$. Đặt
$a=1+|A|+|B|$ và $b=1+|C|+|D|$, ta được
\[
  \lfloor\log a\rfloor
  = \lfloor\log b\rfloor
  = \lfloor\log(a+b+1)\rfloor. \tag{1}
\]
Không mất tính tổng quát, giả sử $b\ge a$. Khi đó
\[
  \lfloor\log(a+b+1)\rfloor
  \ge \lfloor\log(2a)\rfloor
  = 1+\lfloor\log a\rfloor
  > \lfloor\log a\rfloor, \tag{2}
\]
mâu thuẫn với (1). Do đó khi $\mu'(x)=\mu(x)$, bước Zig-Zig không cần chi phí
thêm, và hiển nhiên cũng không vượt quá $3(\mu'(x)-\mu(x))$.

\paragraph{Trường hợp 3: Zig-Zag hoặc Zag-Zig.}
Tương tự trường hợp thứ hai, có thể chứng minh mỗi bước Zig-Zag cũng tốn
không quá $3(\mu'(x)-\mu(x))$.

Ba trường hợp trên cho thấy một bước Zig có chi phí tối đa
$3(\mu(S)-\mu(x))+1$, còn một bước Zig-Zig hoặc Zig-Zag có chi phí tối đa
$3(\mu'(x)-\mu(x))$. Cộng chi phí của các phép quay theo thứ tự thực hiện, ta
thấy một thao tác $\operatorname{Splay}(x,S)$ không tốn quá $3\mu(S)+1$.
Do đó mọi thao tác cơ bản dựa trên splay đều có độ phức tạp khấu hao
$O(\log n)$. Vì vậy cây splay là một cấu trúc dữ liệu có hiệu quả thời gian
rất tốt.

\section{Ứng dụng của cây splay}

Là một cấu trúc dữ liệu có hiệu quả thời gian cao và yêu cầu bộ nhớ nhỏ, cây
splay có nhiều đất dụng võ trong giải bài. Ta xét ví dụ sau.

\subsection*{Ví dụ: Thống kê doanh thu Turnover}

Bài toán xuất phát từ kỳ tuyển chọn đội Hồ Nam năm 2002.

\paragraph{Tóm tắt đề bài.}
Tiger vừa được công ty bổ nhiệm làm trưởng phòng kinh doanh. Nhiệm vụ đầu
tiên của anh là thống kê và phân tích tình hình kinh doanh của công ty từ
ngày thành lập. Sổ sách ghi lại doanh thu của từng ngày. Do ngày lễ, giảm giá
lớn hoặc các nguyên nhân khác, doanh thu có thể dao động; dao động vừa phải
là chấp nhận được, nhưng nếu một ngày doanh thu tăng hoặc giảm đột ngột thì
có thể tình hình kinh doanh đang có vấn đề. Trong quản lý kinh tế, người ta
dùng giá trị dao động nhỏ nhất của một ngày để đo hiện tượng này:
\[
  \text{dao động nhỏ nhất của ngày hiện tại}
  =
  \min\{\,|\text{doanh thu một ngày trước đó}
          -\text{doanh thu ngày hiện tại}|\,\}.
\]
Dao động nhỏ nhất càng lớn thì tình hình kinh doanh càng kém ổn định. Để
phân tích mức ổn định của toàn bộ công ty từ khi thành lập đến nay, chỉ cần
cộng các giá trị dao động nhỏ nhất của từng ngày. Nhiệm vụ là viết chương
trình giúp Tiger tính tổng này.

Ngày đầu tiên có giá trị dao động nhỏ nhất bằng chính doanh thu của ngày đó.
Dữ liệu có $n\le 32767$, doanh thu mỗi ngày $a_i\le 1\,000\,000$, và kết quả
cuối cùng $T\le 2^{31}$.

\paragraph{Phân tích ban đầu.}
Ý nghĩa của đề rất rõ: mỗi lần đọc một số mới, ta cần tìm trong các số đã
đọc trước đó một số có khoảng cách nhỏ nhất với nó.

Ta dễ nghĩ tới thuật toán $O(n^2)$: mỗi lần đọc một số, duyệt lại tất cả các
số trước đó để lấy hiệu nhỏ nhất với số hiện tại rồi cộng vào kết quả $T$.
Nhưng vì $n$ lớn, cách này không thể chạy kịp thời hạn. Nếu dùng cây đoạn để
ghi nhận các giá trị đã xuất hiện, ta cần một mảng cỡ $2M$, điều này không
khả thi trong môi trường Turbo Pascal 7.0 của cuộc thi khi đó. Cây đỏ đen và
cây AVL có thời gian và bộ nhớ tốt, nhưng độ phức tạp lập trình cao. Vì vậy
cây splay trở thành một lựa chọn tự nhiên.

\paragraph{Mô tả thuật toán.}
Phân tích kỹ hơn, bài toán chỉ cần ba thao tác trên tập có thứ tự: chèn, tìm
tiền nhiệm và tìm kế nhiệm. Cây splay có độ phức tạp rất tốt cho cả ba thao
tác này, nên ta thiết kế thuật toán như sau.

Ban đầu cây $S$ rỗng, tổng $T$ bằng $0$. Mỗi lần đọc một số $p$, thực hiện
$\operatorname{Insert}(p,S)$ để chèn $p$ vào cây splay $S$. Sau bước này, $p$
được đưa lên gốc. Khi đó nút phải nhất trong cây con trái của $p$ là tiền
nhiệm của $p$, còn nút trái nhất trong cây con phải của $p$ là kế nhiệm của
$p$ trong tập có thứ tự. Lấy hiệu nhỏ nhất giữa $p$ với hai giá trị này rồi
cộng vào kết quả $T$.

\paragraph{Nhận xét.}
Vì các thao tác cơ bản của cây splay đều có độ phức tạp khấu hao
$O(\log n)$, toàn bộ thuật toán có độ phức tạp thời gian $O(n\log n)$ và có
thể chạy trong thời hạn. Về bộ nhớ, ta có thể dùng mảng để mô phỏng con trỏ
và lưu cấu trúc cây; lượng bộ nhớ không quá $400$ KB, đủ để dùng bộ nhớ động
trong Pascal. Về lập trình, cây splay khá đơn giản, chương trình không quá
phức tạp. Dù không phải thuật toán duy nhất cho bài này, cây splay vẫn có
nhiều ưu thế so với các cấu trúc dữ liệu thường dùng khác.

\begin{center}
\begin{tabular}{lcccc}
\hline
 & Tìm tuần tự & Cây đoạn & Cây AVL & Cây splay \\
\hline
Độ phức tạp thời gian & $O(n^2)$ & $O(n\log a)$ & $O(n\log n)$ & $O(n\log n)$ \\
Độ phức tạp bộ nhớ & $O(n)$ & $O(a)$ & $O(n)$ & $O(n)$ \\
Độ phức tạp lập trình & rất đơn giản & khá đơn giản & khá phức tạp & khá đơn giản \\
\hline
\end{tabular}
\end{center}

Bảng trên cho thấy trong bài này cây splay đều khá tốt ở ba tiêu chí thời
gian, bộ nhớ và độ khó lập trình, nên rất phù hợp để dùng trong thi đấu.

\section{Tổng kết}

Từ phần phân tích trên, ta thấy cây splay có các ưu điểm sau.

\begin{enumerate}
  \item Độ phức tạp thời gian thấp. Mọi thao tác cơ bản của cây splay đều có
    độ phức tạp khấu hao $O(\log n)$, đây là mức rất tốt trong các cấu trúc
    dữ liệu dạng cây.
  \item Yêu cầu bộ nhớ không cao. Khác với cây đỏ đen phải lưu màu của mỗi
    nút, hay cây AVL phải lưu hệ số cân bằng, cây splay không cần lưu thông
    tin phụ để duy trì cân bằng.
  \item Thuật toán đơn giản, dễ lập trình. Các thao tác cơ bản của cây splay
    đều dựa trên thao tác splay; trong thao tác này chỉ cần căn cứ vào vị trí
    của nút hiện tại để thực hiện các phép quay.
\end{enumerate}

Mặc dù cây splay và cây AVL không chênh lệch nhiều về độ phức tạp thời gian,
thậm chí trong một số tình huống cây splay có thể chậm hơn AVL, nhưng độ phức
tạp lập trình của cây splay thấp hơn AVL rất nhiều. Trong thi đấu, sử dụng
cây splay thường thuận lợi hơn khi viết và gỡ lỗi.

Trong thi tin học, không nên chỉ theo đuổi thuật toán có hiệu quả thời gian
cao nhất. Cần tìm một điểm cân bằng giữa độ phức tạp thời gian, độ phức tạp
bộ nhớ và độ phức tạp lập trình, rồi lựa chọn thuật toán hợp lý. Điều này đòi
hỏi ta thường xuyên suy ngẫm và nghiên cứu sâu các thuật toán trong quá trình
luyện tập, để khi thi có thể vận dụng thành thạo.

\section*{Tài liệu tham khảo}

\begin{enumerate}
  \item Fu Qingxiang, Wang Xiaodong. \textit{Algorithms and Data Structures}.
    Publishing House of Electronics Industry, 1998.01.
  \item Yan Weimin, Wu Weimin. \textit{Data Structures}, second edition.
    Tsinghua University Press, 1992.06.
  \item \textit{Introduction to Algorithms}, Second Edition, 2001.
\end{enumerate}

\section*{Phụ lục}

\begin{enumerate}
  \item Cài đặt thao tác splay và các thao tác cơ bản khác được tham chiếu
    trong tài liệu gốc qua tệp \texttt{Splay Tree.doc}.
  \item Quá trình cụ thể của các thao tác cơ bản trên cây splay được minh họa
    trong tài liệu gốc qua hoạt hình \texttt{Splay Tree.htm}.
  \item Với bài ví dụ trong bài viết, tác giả đã viết chương trình
    \texttt{Turnover.pas} sử dụng thuật toán cây splay.
\end{enumerate}

\end{document}
